% =========================================================
% Proof: Stern--Brocot Interval Lengths Shrink to Zero
% =========================================================

\newpage
\phantomsection
\label{prf:sb-length-zero}

\begin{remark*}[Return]
\hyperref[lem:sb-length-zero]{Return to Lemma~\ref*{lem:sb-length-zero}.}
\end{remark*}

\begin{theorem*}[Stern--Brocot interval lengths shrink to zero]
Let $(I_n)_{n \geq 0}$ be the nested rational intervals generated by the
Stern--Brocot refinement process, with $I_n = [a_n/b_n, c_n/d_n]$.
Then $\ell_n = c_n/d_n - a_n/b_n \to 0$.
\end{theorem*}

\begin{proof}[Proof (Professional Standard)]
At every level the two boundary fractions are Farey neighbors (the
determinant condition $b_n d_n' - a_n c_n' = 1$ is preserved by mediant
insertion, as established in the proof of
\hyperref[lem:sb-denom-growth]{Lemma~\ref*{lem:sb-denom-growth}}).
By the \hyperref[thm:farey-gap]{gap between Farey neighbors},
\[
\ell_n = \frac{c_n}{d_n} - \frac{a_n}{b_n} = \frac{1}{b_n d_n}.
\]
By \hyperref[lem:sb-denom-growth]{Lemma~\ref*{lem:sb-denom-growth}},
$b_n \to \infty$ and $d_n \to \infty$.  Therefore $1/(b_n d_n) \to 0$,
i.e.\ $\ell_n \to 0$.
\end{proof}

\begin{proof}[Proof (Detailed Learning)]
\textbf{Step 1: Identify the boundary fractions as Farey neighbors.}
At each level $n$, the Farey Neighbor Theorem applies to the two boundary
fractions $a_n/b_n$ and $c_n/d_n$ because mediant insertion preserves
the determinant condition $b_n c_n - a_n d_n = 1$
(as verified in \hyperref[prf:sb-denom-growth]{the denominator growth proof}).

\textbf{Step 2: Apply the Farey gap formula.}
By the \hyperref[thm:farey-gap]{Gap Between Farey Neighbors}
(\hyperref[thm:farey-gap]{Theorem~\ref*{thm:farey-gap}}):
\[
\ell_n = \frac{c_n}{d_n} - \frac{a_n}{b_n} = \frac{b_n c_n - a_n d_n}{b_n d_n} = \frac{1}{b_n d_n}.
\]

\textbf{Step 3: Show $1/(b_n d_n) \to 0$.}
Let $\varepsilon > 0$.  By the Archimedean property
(\hyperref[thm:archimedean-q]{Theorem~\ref*{thm:archimedean-q}}),
there exists $M \in \mathbb{N}$ with $M > 1/\varepsilon$.  By
\hyperref[lem:sb-denom-growth]{Lemma~\ref*{lem:sb-denom-growth}},
$b_n \to \infty$ and $d_n \to \infty$, so there exists $N \in \mathbb{N}$
such that for all $n \geq N$, $b_n > M$ and $d_n > M$.

\textbf{Step 4: Bound $\ell_n$ for $n \geq N$.}
For $n \geq N$:
\[
\ell_n = \frac{1}{b_n d_n} < \frac{1}{M \cdot M} = \frac{1}{M^2} < \frac{1}{M} < \varepsilon.
\]

\textbf{Step 5: Conclude.}
Since $\varepsilon > 0$ was arbitrary, $\ell_n \to 0$.
\end{proof}

\begin{remark*}[Proof structure]
The proof chains two results: the Farey gap formula gives the exact
expression $\ell_n = 1/(b_n d_n)$, and the Denominator Growth Lemma
gives $b_n, d_n \to \infty$.  The Archimedean property converts the
asymptotic statement into a concrete $N$.
\end{remark*}

\begin{remark*}[Dependencies]
\hyperref[thm:farey-gap]{Gap between Farey neighbors},
\hyperref[lem:sb-denom-growth]{Denominator growth on infinite SB paths},
\hyperref[thm:archimedean-q]{Archimedean property of $\mathbb{Q}$}.
\end{remark*}
